\chapter*{A note on the interface language (i18n)}
\addcontentsline{toc}{chapter}{A note on the interface language (i18n)}
\markboth{A note on the interface language}{A note on the interface language}

MobileCAD 0.5.18 is the first build whose interface can be Chinese: it follows the device language, and
\menupath{Preferences → Language} pins English or Simplified Chinese. The screenshots in this manual show the
Chinese build, because the manual was written together with that localisation.

This page explains the mechanism briefly, so the manual can also serve as an i18n walkthrough.

\section*{How the language is chosen}
\begin{itemize}
  \item On start, \file{I18n.init(context)} reads the saved choice (\ui{system}, \ui{en} or \ui{zh-Hans}); with
        \ui{system} it looks at \file{Locale.getDefault().language} and loads \file{assets/i18n/zh-Hans.tsv} for
        Chinese, or keeps English for everything else.
  \item Changing the choice writes a preference and recreates the activity; the document and the engine session
        survive, so the switch is instant and lossless.
\end{itemize}

\section*{The catalog}
\begin{itemize}
  \item One line per string: \file{English<TAB>translation}, UTF-8, `\#` starts a comment. \textbf{The English text
        is the key}; a string the catalog does not contain stays English, so a partial catalog can never break the
        interface.
  \item Translation happens where the text is drawn: the command catalog is translated once when it is loaded
        (the English label stays available for tool search), menus, dialogs, banners, messages, task panels,
        property labels and the NaviCube faces are translated at render time. Document, object and layer names —
        user data — are never looked up.
  \item Commands, task-panel fields and property names use FreeCAD's own Chinese wording. That part of the catalog
        is generated from FreeCAD's \file{src/**/translations/*\_zh-CN.ts} with
        \file{python3 tools/i18n\_import\_freecad.py <FreeCAD checkout>} (LGPL-2.1-or-later), so both applications
        speak the same Chinese.
  \item Adding a language is one more \file{assets/i18n/<code>.tsv} plus its code in \file{I18n.choices}; the
        fallback and the draw-time lookups stay unchanged.
\end{itemize}

\section*{Testing a translation quickly}
\begin{enumerate}
  \item Set the device language (or \menupath{Preferences → Language}) to the target language and restart.
  \item Walk through: menus, a tool panel (command names and tooltips), a feature dialog (property labels), a
        FreeCAD task panel (BIM/FEM/CAM), the constraint list, the language picker itself and the NaviCube.
  \item Switch to English and back: the interface must redraw and the open document must survive.
  \item Load a file made before the change: geometry must be identical; only text changed.
\end{enumerate}

\section*{What stays English}
Engine-side strings (Python exceptions, some FreeCAD C++ messages), the long \ui{About} text and the license
screen are not translated yet. They are listed in the manual's FAQ and are good first contributions.
